% =====================================================================
% conclusion.tex — Section VI: Conclusion and Future Work
%
% Usage: included from main.tex via \input after Section V (Analysis and
% Discussion). No extra packages required; compatible with IEEEtran.
%
% \ref targets below are wired to the real labels in the other sections:
%   \ref{subsec:rq1}             -> Section V-A   (RQ1: Structured Graph Retrieval)
%   \ref{subsec:rq2}             -> Section V-B   (RQ2: Ranked retrieval mechanism)
%   \ref{subsec:rq3}             -> Section V-C   (RQ3: Comparison against others)
%   \ref{subsec:limitations}     -> Section V-D   (Limitations)
%   \ref{subsubsec:system-level} -> Section III   (System-Level Evaluation design)
%   \ref{sec:results-system}     -> Section IV    (System-Level Evaluation results)
%   \ref{sec:results-quant}      -> Section IV    (Embedding Quantisation results)
%   \ref{sec:lr-graph}           -> Section II    (Graph-augmented retrieval)
%
% Search for "[FILL" for the RQ3 findings slot, and for "% NOTE" for
% author decisions to confirm before the final submission.
%
% Quantisation-results pass (July 2026): one sentence appended to the RQ2
% conclusion, the memory-budget clause in the partner-outcomes paragraph
% extended with the priced operating points, and a new future-work item
% (rank agreement at production scale) inserted after the Boolean-gate item.
% =====================================================================

\section{Conclusion and Future Work}
\label{sec:conclusion}

This thesis set out to reconcile two retrieval paradigms that deployed literature-search engines force a choice between: structured retrieval, which answers a query deterministically and auditably, and semantic retrieval, which ranks papers by meaning and so can surface relevant work whose wording differs from the query. It did so in an industry setting, designing, building, and evaluating a production-ready scholarly search system in which a typed knowledge graph of approximately 150 million biomedical works, ingested from OpenAlex, anchors a three-stage cascade of structured graph retrieval, dense embedding retrieval, and cross-encoder re-ranking. Three research questions structured the evaluation, one per layer of the design. This section states what the results and analysis establish for each, summarises the outcomes of the work for the industry partner, and identifies the work that follows from its limitations.

\subsection{Conclusions}
\label{sec:conclusion-conclusions}

In answer to RQ1, modern language models can construct both artefacts that deterministic structured graph retrieval requires, though to differing degrees. On the query side, the selected model translates natural-language queries into the intermediate representation at a joint accuracy of 0.937 on the hand-built gold set, with entity extraction near ceiling and the Boolean expression the binding component. Its residual failures concentrate in the structurally hardest strata and in a categorisation slip that tighter prompting can address (Section~\ref{subsec:rq1}). On the corpus side, the same model produces the strongest benchmarked keyword index of any candidate across three scientific keyphrase datasets, in a task whose absolute scores are constrained by well-documented ceilings in the metric (Section~\ref{sec:lr-keyphrase}). Both artefacts carry the characteristic cost of construction by a language model, in that neither stage is fully reproducible across runs, and the small generated share the keyword prompt intends is not guaranteed faithful to the source in the way a purely extractive method is. Because both artefacts feed a deterministic matcher, the quality of the structured stage as a whole is bounded by how well they can be built, and the evidence places that bound high on the query side and competitive but imperfect on the corpus side.

In answer to RQ2, current bi-encoders and cross-encoders perform scientific ad-hoc search well enough to equip both semantic stages of the cascade, with the leaders of each stage clustered at averages around 82 to 83 nDCG and their confidence intervals overlapping, so the benchmarks do not separate them.
% CHANGED (item 2, bootstrap audit): was "separated by less than their
% confidence intervals", which reads overlap as a positive finding of
% equality. Reworded to the conservative reading defined in Section III-D. Two findings shaped the selection more than raw accuracy. First, the models that lead are recent general-purpose families rather than domain- or task-specific ones. Specialisation helped only in-domain, and the biomedical reranker did not lead even on biomedical text. Second, the paid hosted models do not significantly outperform the best open models in either stage, so the choice legitimately turned on deployment characteristics, namely the bi-encoder's permanent memory cost, which scales with embedding dimension across 150 million stored vectors, against the cross-encoder's per-query latency, which scales with parameter count. Assigning both stages to hosted models is therefore an engineering decision rather than a quality one, and a reversible one, since accuracy-viable open alternatives were identified for both stages (Section~\ref{subsec:rq2}). The storage ablation closes the remaining open cost of the dense stage. Quantising the stored vectors to int8 is measurably free on the benchmarks, and the deployed embedder retains 99.8 percent of its average benchmark quality at half its stored footprint, pending validation at production pool sizes (Section~\ref{sec:results-quant}).

RQ3 asks how the assembled cascade's results differ from those of existing search systems and what explains the differences. The component results ground clear expectations. The cascade should be strongest where queries carry typed structure, such as authors, institutions, venues, and Boolean composition, which neither a semantic-only retrieval-augmented LLM nor Google Scholar's opaque ranking expresses; it should be precision-oriented by construction, returning only works that satisfy the Boolean gate; and it should be bounded where its substrate is, silently excluding works whose keywords or affiliations the index lacks. The system-level evaluation specified in Section~\ref{subsubsec:system-level} tests these expectations by comparing the cascade against a retrieval-augmented LLM and Google Scholar over a curated set of realistic queries, characterising per query what each system surfaces and misses under light relevance judgments, and its results are reported in Section~\ref{sec:results-system} and analysed in Section~\ref{subsec:rq3}. \emph{[FILL once the evaluation has run: two to four sentences stating the observed pattern (where the structured gate gave the cascade results the comparators could not reach, where it cost recall relative to them, one representative query in each direction) and the verdict on whether the expectations above held.]} One limit applies to whatever the comparison finds, and is stated here so that the conclusion carries it. The instrument is qualitative by design. It characterises behaviour and mechanism rather than measuring effect sizes, because no ground-truth relevance pool exists over a live 150-million-work corpus, and light judgments over a small query set cannot support statistically powered claims. Its findings are therefore observational, and converting them into quantitative ones is the first item of future work below.

Taken together, the component evidence supports the design claim of this thesis, that structured and semantic retrieval need not be traded off. A language model can be confined to the two tasks it is demonstrably good at, building the keyword index offline and translating query intent online, while deterministic compilation preserves the correctness and auditability that systematic search requires, and established dense models supply the semantic ranking that exact matching cannot. What the component scores cannot establish alone is how the composed system behaves against the tools researchers already use, and that confirmation rests on the system-level evaluation. The thesis's contributions stand independently of that confirmation: the typed graph and its documented ingestion pipeline; the natural-language query interface that decouples intent extraction from query compilation; the stratified translation benchmark that no existing benchmark covered; and the component evaluations that make the accuracy and deployment trade-offs of every model choice explicit.

\subsection{Outcomes for the Industry Partner}
\label{sec:conclusion-partner}

The most direct outcome is the system itself. The knowledge graph, ingestion pipeline, and retrieval cascade this thesis documents form a production-ready system for Latent Knowledge to deploy, and every learned component in that pipeline runs the model that the benchmarks reported here selected. The production configuration consequently rests on measured evidence rather than default choice: the selected keyword extractor is the strongest of the candidates benchmarked, the selected translation model tops its benchmark, and the selected embedder and reranker sit inside the leading statistical cluster of their fields.
% NOTE: Table XV marks text-embedding-3-small (and SPECTER2-base) as
% "(existing)". If that means the embedder was already the incumbent in
% production before this work, strengthen the sentence above to say the
% benchmark *validated the incumbent against the field* rather than
% assumed it --- a stronger partner outcome. Confirm before final.

The evaluations also leave the partner with decision assets beyond the selections themselves. Because the accuracy race in both semantic stages ended in statistical ties, the analysis could isolate the deployment characteristics that actually separate the candidates: the memory budget each embedding dimension implies over the corpus, around 115\,GB at 768 dimensions against roughly 230\,GB at 1536, now priced by the quantisation matrix down to 99.8 percent of benchmark quality at half the footprint without re-encoding; the licence terms that exclude particular open models from commercial use; and the axis each stage's cost scales on as the user base grows. The practical consequence is a costed and reversible path off the paid APIs, with EmbeddingGemma-300M and Qwen3-Reranker-4B identified as accuracy-viable replacements should scale or cost structure make self-hosting preferable, and with the one hard dependence, hosted IR extraction, named explicitly rather than discovered later.

Finally, the evaluation apparatus outlives the choices it made. The hand-built stratified gold set and the component harnesses together form a regression suite for a pipeline whose learned components are served by third parties and can change without notice, and the limitations gathered in Section~\ref{subsec:limitations} function as a risk register for the deployment, in which the substrate's metadata gaps, the reproducibility cost of artefacts built by language models, and the provider dependence are each named, sized where the evidence allows, and paired with a mitigation.

\subsection{Future Work}
\label{sec:conclusion-future}

Each item below is the implication of a limitation identified in Section~\ref{subsec:limitations} or of an open thread in the analysis, ordered from nearest term to most exploratory.

\begin{itemize}
  % NOTE (July 2026 edit pass): the "complete the evaluation" half of this item was
  % removed because the Section III-D6 protocol will run before final submission. If the
  % pooled two-annotator quantitative protocol will ALSO run before submission, delete
  % this item entirely and update the forward reference at the end of the RQ3 paragraph
  % in VI-A ("the first item of future work below").
  \item \textbf{Quantify the system-level evaluation.} Pooled top-$k$ relevance judgments from at least two annotators with agreement reported, and graded metrics computed over the pooled results, would turn the observational findings about what each system surfaces into measured ones.

  \item \textbf{Measure the recall cost of the Boolean gate.} Stage-one recall for topical entities is bounded by the offline keyword index, and no later stage can recover a work the gate excludes. Relaxing the gate on a sampled query set, re-ranking the widened pool, and counting the relevant works admitted would estimate the practical size of a bound the component benchmarks cannot see.

  \item \textbf{Validate reduced-dimension operating points at production scale.} The retention measured by the quantisation matrix is conditioned on benchmark pool sizes (Section~\ref{subsec:limitations}). Measuring top-100 overlap between the native and a reduced configuration over candidate pools sampled from the live graph would test whether that retention transfers to the pools the cascade actually ranks, needs no relevance labels, and answers the operational question of whether the candidates handed to the cross-encoder change. Binary quantisation is the natural next point on the same frontier.

  \item \textbf{Quantify and surface substrate gaps.} The affiliation gap documented for OpenAlex should be measured on the biomedical subset specifically, since the published figures are corpus-wide, and the retrieval path should surface the exclusion to the user, for example by reporting how many works were removed from an institution-anchored query for missing affiliations, so that the precision-oriented behaviour is transparent rather than silent. A complementary refinement models the affiliation's time context rather than only its presence. Because a Person's Institution changes over a career, an association node between the two carrying the period of the affiliation would let institution-anchored queries respect the year a work was actually written rather than treating affiliation as fixed.

  \item \textbf{Harden the translation gold set.} A second annotator with inter-annotator agreement reported would address the single-author construct-validity risk; sourcing queries from production logs as the user base grows would ground the set in real usage rather than designer expectation; and the hardest strata, where per-stratum counts are smallest, should be expanded first.

  \item \textbf{Unwind the hosted dependence for IR extraction.} The gap between the hosted and open models is widest on this task, at 0.937 against 0.802 for the best open model, which makes it the dependence hardest to reverse. Fine-tuning an open model on an expanded gold set, or distilling from the hosted model, is the direct path. In the nearer term, tighter prompting that defines filter terms explicitly targets the largest single error source.

  \item \textbf{Exploit the graph beyond candidate selection, and test the schema's generality.} The typed graph and Memgraph's algorithm library make citation-graph signals, such as Personalized PageRank seeded from the matched entities, available as an additional ranking signal, which is the hybridisation the literature identifies as conceptually well-matched but unbenchmarked (Section~\ref{sec:lr-graph}). Ingesting a second, non-biomedical field would convert the schema's claimed generality from a design property into a measured one.
\end{itemize}

Models better suited to every learned task in this pipeline will continue to emerge. The architecture was designed so that adopting them is a substitution rather than a redesign, and that property, more than any single selection reported here, is what this work leaves the deployment with.